\documentclass[journal]{IEEEtran}

\usepackage{cite}
\usepackage{amsmath,amssymb,amsfonts}
\usepackage{graphicx}
\usepackage{textcomp}
\usepackage{xcolor}
\usepackage{geometry}
\usepackage{svg}

\import\texenv

\newcommand{\nwcv}[2]{\overline{#1#2^*}}
\newcommand{\nwv}[1]{\overline{|#1|^2}}

\setlength{\parskip}{0.5em}
\setlength{\parindent}{0pt}

\geometry{portrait, top=0.8in, bottom=0.5in, left=0.7in, right=0.7in, includefoot}

\def\BibTeX{{\rm B\kern-.05em{\sc i\kern-.025em b}\kern-.08em
    T\kern-.1667em\lower.7ex\hbox{E}\kern-.125emX}}
	
\begin{document}

\title{Noise Wave Analysis Techniques for Multi-port Networks}
\author{ECE 6323,  Rick Lyon}
\maketitle
\begin{abstract}
This paper is based on a publication in IEEE Transactions On Microwave Theory and Techniques that outlines an analytical method for calculating the noise performance of multi-port RF networks.\cite{b3} The technique is implemented in C++ and compared to the commercially available circuit solver ADS.
\end{abstract}

\section{Introduction}

The noise figure of a  two port receiver chain can be estimated by cascading the individual noise figures and gain from each component, but this approach does not work for more complicated like the one shown in Figure \ref{fig:img/parallel}. Noise wave analysis models noise as scattered waves generated by normal random variables and can accurately predict noise performance of complicated multi-port networks.

\texenv\figure[file=img/parallel, caption={\small Networks with amplifiers coupled at the input require advanced noise analysis techniques. For high gain amplifiers, even a small amount of coupling at the input can lead to significant deviations from the estimated noise figure using traditional methods.}, width=3.2in]

\section{Noise Wave Theory}


The time-averaged voltage output of an RF noise source is zero, but the average power is non-zero and can be quantified with the root mean squared voltage (RMS). The RMS value of a waveform \(v(t)\) is defined as,
\[v_\textrm{rms} = \sqrt {{1 \over {T}}{\int _{0}^{T}{v(t)}^{2}\,dt}}\]
For sinusoidal waveforms, this reduces to \(v(t) / \sqrt{2}\), but waveforms produced by noise sources are not sinusoidal and the full RMS expression is used. The distribution of voltages produced by a noise source can also be represented with a random variable \(V\) with variance,
\[ var(V) = \mathbf{E}[V^2] \] 

The expectation of \(V^2\) averages the square of the samples produced by \(V\). The expectation can also be found by averaging the voltage waveform in the usual way,
\begin{equation}\mathbf{E}[V^2] = \lim_{T\to\infty}{1 \over {T}}{\int _{0}^{T}{v(t)}^{2}\,dt} \label{eq:rms}\end{equation}

If we assume \(T\) is sufficiently large, 
\[v_\textrm{rms}^2 = var(V)\]
Which means that the standard deviation \(\sigma\) of the random variable \(V\) is equal to the RMS voltage,
\[\sigma = \sqrt{var(V)} = v_\textrm{rms}\]

Noise wave theory uses this result to model noise produced by RF devices with a random variable \(c_n\) at each port. This represents an outgoing noise waveform with a 1 Hz bandwidth, equivalent to the noise spectral density. The RMS voltage of \(c_n\) as shown in Figure \ref{fig:img/cmatrix} is,
\[v_\textrm{rms}^2 = var(c_n) = \mathbf{E}[c_n c_n^*] = \overline{|c_n|^2}\]

\texenv\figure[file=img/cmatrix, caption={\small Noise generated by an RF device is modeled as an outgoing wave from each port created by the normal random variable \(c_n\)}., width=3.5in, placement=b!]

The overbar denotes time averaging as shown in (\ref{eq:rms}), and the expectation was generalized for complex random variables. Noise waves propagate through the network and will scatter just like normal voltage waves, incurring phase shifts and amplitude scaling, which necessitates the use of complex random variables. The difference from normal voltage waves however, is that they do not always combine with other sources coherently.


\subsection{Adding Noise Waves}

The sum of two random variables \(V_1\) and \(V_2\) is \cite{b1},
\[var(V_1+V_2) = var(V_1) + var(V_2) + 2cov(V_1,V_2)\]

In the case of noise waves propagating through a network, the voltage waves \(c_n\) are complex variables and may be scaled by some arbitrary complex number \(a\) or \(b\). It helps to derive the above result in this more general case,
\begin{gather}
var(ac_1+bc_2) = \mathbf{E}\Big[(ac_1+bc_2)(ac_1+bc_2)^*\Big] \notag \\
= \mathbf{E}\Big[|a|^2|c_1|^2+ba^*c_2c_1^*+ab^*c_1c_2^*+|b|^2|c_2|^2\Big] \label{eq:var_sum}
\end{gather}
%
Because the expectation is a time-averaged integration as shown in (\ref{eq:rms}), the expectation operator is distributive.
%
\begin{gather}
var(ac_1+bc_2) = \notag\\
|a|^2\overline{|c_1|^2}+ba^*\overline{c_2c_1^*}+ab^*\overline{c_1c_2^*}+|b|^2\overline{|c_2|^2} \label{eq:var_sum1}
\end{gather}
%

If \(c_1\), \(c_2\), \(a\) and \(b\) are real rather than complex, this reduces to the more familiar result,
\begin{gather}
var(ac_1+bc_2) = \notag\\
|a|^2 var(c_1) + |b|^2 var(c_2) +  (2ab)cov(c_1, c_2)
\end{gather}
%

The four terms in (\ref{eq:var_sum1}) can be arranged into a correlation matrix. For a two-port device, the noise waves at each port is characterized with,
\[ \mathbf{C_s} =
\begin{pmatrix}
\overline{|c_1|^2} & \overline{c_1c_2^*} \\
\overline{c_2c_1^*} & \overline{|c_2|^2}
\end{pmatrix}
\]
%
Because \(\overline{|c_1|^2} = v_\textrm{rms}^2\), the diagonal terms are the impedance-normalized noise powers generated at the output of each port. The off-diagonal terms are the covariances between each port. This generalizes for n-port devices since the result in (\ref{eq:var_sum}) is easily extended to three or more random variables. A N-port device will have a NxN correlation matrix.

The total noise output power from a device is a combination of the internally generated noise represented by \(c_n\), and the input noise at each port. If the input noise is represented by random variables \(a_1\) and \(a_2\), the resulting random variables at the device ports \(b_1\) and \(b_2\) are,

\begin{gather}
	\begin{bmatrix}
	b_1 \\
	b_2 
	\end{bmatrix}=
	\begin{bmatrix}
	S_{11} & S_{12} \\
	S_{21} & S_{22}
	\end{bmatrix}
	\begin{bmatrix}
	a_1 \\
	a_2 
	\end{bmatrix} +
	\begin{bmatrix}
	c_1 \\
	c_2 
	\end{bmatrix}
	\label{eq:b_matrix}
\end{gather}

The variances (noise power) of \(b_1\) and \(b_2\) can be computed using the correlation matrix terms of \(\mathbf{C_s}\) and the result shown in (\ref{eq:var_sum}). The covariance between the output noise from each port is,
%
\begin{gather}
cov(b_1, b_2) = \notag \\
\mathbf{E}\Big[(S_{11}a_1 + S_{12}a_2 + c_1)(S_{21}a_1 + S_{22}a_2 + c_2)^*\Big] \notag\\
= S_{11}S_{21}^*\overline{|a_1|^2} + S_{12}S_{22}^*\overline{|a_2|^2} + \overline{c_1c_2^*} \label{eq:cov}
\end{gather}
%
Where it is assumed that \(a_1\) and \(a_2\) are uncorrelated with each other, and uncorrelated with the internally generated noise variables \(c_1\) and \(c_2\). For example, \(\overline{a_1c_1^*}= 0\). The result in (\ref{eq:cov}) is interesting because even though the internally generated noise sources are uncorrelated, the output waves at the ports are correlated since scaled versions of \(a_1\) and \(a_2\) combine at both ports.

The results in this section can be extended to determine the aggregate correlation matrix of a network of interconnected devices, but first the correlation matrix of each device in the network must be known. The next two sections describe how to derive these matrices for passive and active devices.

\subsection{Noise Wave Terms for Passive Devices}
If a passive component is in thermal equilibrium with its surroundings, the total noise power leaving and exiting the device must be equal.Otherwise, the net flow of power would lead to a change of temperature in the device. Passive devices are not always in thermal equilibrium with their terminations, but the internally generated noise depends only on the device temperature, and not on the input noise waves. This means that the device noise correlation matrix derived from this special case will hold even when the device is not in thermal equilibrium.

In the case of a two port device, the total input noise power into both ports is \(|a_1|^2 + |a_2|^2 = 2kT\). To maintain thermal equilibrium, the total noise power exiting the device must the same, 
\[\overline{|b_1|^2} + \overline{|b_2|^2} = 2kT\]
%
Since the variance of the sum of the random variables \(b_1\) and \(b_2\) is the total power leaving the device, \cite{b6}
\[var(b_1 + b_2) = 2kT\]
%
Comparing this with the variance shows that \(\overline{b_1 b_2^*} = \overline{b_2 b_1^*} = 0\),
\[\overline{|b_1|^2} + \overline{|b_2|^2} + \overline{b_1 b_2^*} + \overline{b_2 b_1^*} = 2kT\]
%

With the result that, \(\overline{b_1 b_2^*} = 0\), (\ref{eq:cov}) can be solved for the covariance between \(c_1\) and \(c_2\),
\begin{gather}
\overline{c_1c_2^*} = kT\left( -S_{11}S_{21}^* - S_{12}S_{22}^*\right) \label{eq:ccvar}
\end{gather}

The variances of \(c_n\) can be found by computing the variance of \(b_n\) from (\ref{eq:b_matrix}) and solving for \(\overline{|c_n|^2}\),
\begin{gather}
	\overline{|c_1|^2} = kT\Big(1- |S_{11}|^2 - |S_{12}|^2\Big) \notag \\
	\overline{|c_2|^2} = kT\Big(1- |S_{21}|^2 - |S_{22}|^2\Big) \label{eq:cvar}
\end{gather}

This shows that all the noise terms of \(\mathbf{C_s}\) for passive devices can be computed solely from the device s-parameters. (\ref{eq:ccvar}) and (\ref{eq:cvar}) can be conveniently combined into matrix form by inspecting each term,
\[\mathbf{C_s} = kT ( \mathbf{I} - \mathbf{S}\mathbf{S}^H ) \]
Where \(\mathbf{S}\) is the s-matrix, \(\mathbf{I}\) is the identity matrix, and \(^H\) denotes the Hermitian transpose. This equation can be verified for passive s-matrices with any number of ports by solving the variances and covariances in matrix form as done in \cite{b7}. This result is known as Bosma's Theorem. \cite{b6}


\subsection{Noise Wave Terms for Active Devices}
This section derives the noise correlation matrix for active 2-port devices from the device noise parameters. Most of the available literature focuses on novel ways to directly measure the correlation matrix, but in practice, it's more efficient to compute \(\mathbf{C_s}\) from the noise parameters since these are typically available from the manufacturer. The disadvantages of this approach is that noise parameters are only defined for 2-port devices, and noise parameters are typically not as accurate as other methods, but it's sufficient for many practical applications.

We start by considering the available noise power at the output of an active device presented with an input reflection coefficient \(\Gamma_s \neq 0\). As shown in Figure \ref{fig:img/noise_param}, the total noise power is a combination of scattered waves from the random variables \(c_1\), \(c_2\), \(a_1\) and \(a_2\).
\texenv\figure[file=img/noise_param, caption={\small Output noise from an RF amplifier when presented with a input reflection \(\Gamma_s\)}., width=3in, placement=b]
%
\[\Gamma_{out} = S_{22} + {S_{12}S_{21} \Gamma_s \over 1 - S_{11} \Gamma_s} \]
%
\[b_2 = a_1 S_{21} + a_2 \Gamma_{out} + c_2 + c_1\Gamma_s S_{21}\]

The input noise waves \(a_1\) and \(a_2\) are assumed to be uncorrelated, which is not always the case in an arbitrary network of devices. For example, \(a_1\) could contain some component of \(b_2\) if the amplifier has an external feedback path. However, when noise parameters are measured, a matched termination is placed on the output port that prevents any appreciable external feedback path. This means that \(a_1\) and \(a_2\) are effectively uncorrelated. The same is not true of \(c_1\) and \(c_2\). The variance of \(b_2\) is then,
\begin{gather}
\overline{|b_2|^2} = |S_{21}|^2\overline{|a_1|^2} + |\Gamma_{out}|^2\overline{|a_2|^2} + \overline{|c_2|^2} + |\Gamma_s S_{21}|^2 \overline{|c_1|^2} \notag \\
 + 2\textrm{Re}\Big(\Gamma_s S_{21} \overline{c_1c_2^*} \Big) \notag
\end{gather}

The equivalent temperature of a 2-port device is the temperature of a resistor that if placed at the input terminals of noise-free version of the device, would generate the same output power as the noisy device in the absence of any other other noise source. In absence of \(a_1\) and \(a_2\), the output noise power is,
\[N_o = \overline{|c_2|^2} + |\Gamma_s S_{21}|^2 \overline{|c_1|^2} + 2\textrm{Re}\Big(\Gamma_s S_{21} \overline{c_1c_2^*} \Big)\]
If the device was noise-free, the same output power could be obtained if a noisy resistor were placed at the input with temperature,
\begin{gather}
T_e = {N_o \over k|S_{21}|^2} \label{eq:f2temp0}
\end{gather}

The noise temperature can be determined with the measured noise parameters (\(F_{min}\), \(R_N\) and \(\Gamma_{opt}\)) of the device using
\[F = F_{min} + \frac{4R_N}{Z_0} \frac{|\Gamma_S - \Gamma_{opt}|^2}{(1-|\Gamma_S|^2)|1+\Gamma_{opt}|^2}\]
For a two-port device, the temperature is related to the noise figure \(F\) by,
\begin{gather}
F = {T_e \over T_0} + 1 \label{eq:f2temp}
\end{gather}
Converting to temperature,
\begin{gather}
{T_{e} } = {T_{min} } + \frac{4R_NT_0}{Z_0} \frac{|\Gamma_S - \Gamma_{opt}|^2}{(1-|\Gamma_S|^2)|1+\Gamma_{opt}|^2} \label{eq:te_nw}
\end{gather}
Equating the two device temperatures derived from noise wave analysis (\ref{eq:f2temp0}) and the noise parameters (\ref{eq:te_nw}),
\begin{gather}
N_o = k|S_{21}|^2 \Big(T_{min} + t \frac{|\Gamma_S - \Gamma_{opt}|^2}{(1-|\Gamma_S|^2)|1+\Gamma_{opt}|^2} \Big) \label{eq:nout}
\end{gather}
Where \(t\) was used to simplify the expression,
\[t = {4R_NT_0 \over Z_0}\]

The noise wave terms can be found by choosing different values for \(\Gamma_s\) that remove unknowns from the equation. The simplest choice is to set \(\Gamma_s=0\) which removes all but \(\overline{|c_2|^2}\) from \(N_o\) in (\ref{eq:f2temp0}).
\[\overline{|c_2|^2} = k|S_{21}|^2 \Big({T_{min} } + t \frac{|\Gamma_{opt}|^2}{|1+\Gamma_{opt}|^2} \Big)\]

I believe the other terms are calculated by selecting other values of \(\Gamma_S\), although I wasn't able to completely redrive them. \cite{b3} lists the conversion equations, but with no explanations,
\begin{align}
	\overline{|c_2|^2} &= k|S_{21}|^2 \Big({T_{min} } + t \frac{|\Gamma_{opt}|^2}{|1+\Gamma_{opt}|^2} \Big)\label{eq:np1} \\
	\overline{|c_1|^2} &= kT_{min} \Big(|S_{11}|^2 -1 \Big) + kt{|1 - S_{11} \Gamma_{opt}|^2 \over |1 + \Gamma_{opt}|^2} \\
	\overline{c_1 c_2^*} &= { -S_{21}^* \Gamma_{opt}^* kt \over |1 + \Gamma_{opt}|^2} + {S_{11} \over S_{21}} \overline{|c_2|^2} \label{eq:np3}
\end{align}


\subsection{Aggregate Noise Wave Parameters}
The goal of this section is to derive the aggregate noise wave parameters of a network with multiple interconnected devices. We start with a general case of creating a 5-port network by connecting two 3-port devices, as shown in Figure \ref{fig:img/interconnect}. This scenario will then be extended to connections between arbitrary devices. 

\texenv\figure[file=img/interconnect, caption={\small Aggregate 5-port network created from two 3-port devices (D1 and D2)  connected at ports 2 and 1 respectively.}, width=3.5in, placement=b]

After the connection is made, the noise waves \(c_n\) no longer describe the noise power at each port because \(c_n\) combines with scattered noise waves from the other device. The internally generated noise power leaving each port can be found by solving for \(b_n\), when all other external noise sources are set to zero, \(a_n = 0\). Physically, there are no noiseless terminations that prevent any reflections back into the network, but this is used purely as mathematical tool. This reduces the six equations for \(b_n\) to,
\begin{gather}
b_1^1 = b_1^2 S_{12}^1 + c_1^1 \notag\\
b_2^1 = b_1^2 S_{22}^1 + c_2^1 \notag\\
b_3^1 = b_1^2 S_{32}^1 + c_3^1 \notag\\
b_1^2 = b_2^1 S_{11}^2 + c_1^2 \notag\\
b_2^2 = b_2^1 S_{21}^2 + c_2^2 \notag\\
b_3^2 = b_2^1 S_{31}^2 + c_3^2 \notag
\end{gather}
Where the constraint imposed by the port connection was used,
\begin{gather}
a_2^1 = b_1^2 \notag \\
a_1^2 = b_2^1 \notag
\end{gather}
%
The equations for \(b_n\) can be written in matrix form as,
%
\[
\renewcommand*{\arraystretch}{1.2}
\begin{bmatrix}
c_1^1 \\
c_2^1 \\
c_3^1 \\
c_1^2 \\
c_2^2 \\
c_3^2 \\
\end{bmatrix}
=
\begin{bmatrix}
1            & 0         & 0   & -S_{12}^1     & 0  & 0 \\
0            & 1         & 0   & -S_{22}^1     & 0  & 0 \\
0            & 0         & 1   & -S_{32}^1     & 0  & 0 \\
0            & -S_{11}^2  & 0   & 1     & 0  &0\\
0            & -S_{21}^2  & 0   & 0     & 1  &0         \\
0            & -S_{31}^2  & 0   & 0     & 0  &1         \\
\end{bmatrix}
\begin{bmatrix}
b_1^1 \\
b_2^1 \\
b_3^1 \\
b_1^2 \\
b_2^2 \\
b_3^2 \\
\end{bmatrix}
\]
Or, written more compactly,
\[\mathbf{c}_t = \mathbf{Q} \mathbf{b}_t \]

The matrix \(\mathbf{Q}\) has a very predictable structure for arbitrary port connections between two devices of any size. A computer program can easily compile \(\mathbf{Q}\) by taking the \(i^{th}\) s-matrix column of the device being connected at port \(i\), and inserting it into an identity matrix at column \(j\) determined by the port being connected of the other device. This process can be repeated for more than one port connection. Once \(\mathbf{Q}\) is compiled, the outgoing noise waves can be solved for,
%
\[\mathbf{b}_t = \mathbf{A} \mathbf{c}_t \]
Where \(\mathbf{A} = \mathbf{Q}^{-1}\).
The final task is to compute \(\overline{|b_n|^2}\) for each element in the \(b_n\) vector, as well as the covariances between each port. This will form a NxN matrix and will be equivalent to the correlation matrix of the aggregate network. There will be an extra dimension in the correlation matrix for the internally connected port, which can be removed. Although sometimes this internal noise wave is useful and can be used as a ``noise probe" inside of the network.

It helps to drop the matrix notation and look at a single \(b\) when evaluating the variance. We assume here that \(\mathbf{A}\) is a 4x4 matrix, meaning two 2-port devices were connected together. Then the noise wave exiting the \(n^{th}\) port of the aggregate network would be given by,
\[b_n = c_1 A_{n1} + c_2 A_{n2} + c_3 A_{n3} + c_4 A_{n4} \]
The variance of \(b_n\) is found by utilizing the results from (\ref{eq:var_sum1}), and using the fact that covariances between noise waves of different devices are zero, (i.e \(\nwcv{c_{1}}{c_{3}} = 0\))
\begin{gather}
\nwv{b_n} = \notag \\
|A_{n1}|^2 \nwv{c_1} + |A_{n2}|^2 \nwv{c_2} + |A_{n3}|^2 \nwv{c_3} + |A_{n4}|^2 \nwv{c_4} \notag \\
+ A_{n1}A_{n2}^* \nwcv{c_{1}}{c_{2}} + A_{n2}A_{n1}^* \nwcv{c_{2}}{c_{1}} \notag
+ A_{n3}A_{n4}^* \nwcv{c_{3}}{c_{4}} + A_{n4}A_{n3}^* \nwcv{c_{4}}{c_{3}} \notag
\end{gather}
%
The covariance between ports \(j\) and \(k\) is,
\begin{gather}
\nwcv{b_j}{b_k} = \notag \\
A_{j1}A_{k1}^* \nwv{c_1} + A_{j2}A_{k2}^*\nwv{c_2} + A_{j3}A_{k3}^* \nwv{c_3} + A_{j4}A_{k4}^* \nwv{c_4} \notag \\
+ A_{j1}A_{k2}^* \nwcv{c_{1}}{c_{2}} + A_{j2}A_{k1}^* \nwcv{c_{2}}{c_{1}} \notag
+ A_{j3}A_{k4}^* \nwcv{c_{3}}{c_{4}} + A_{j4}A_{k3}^* \nwcv{c_{4}}{c_{3}} \notag
\end{gather}
The variance of the noise wave at port \(n\) can be treated as a special case of the covariance equation where \(j=k\). This can be generalized for devices with any number of ports by first compiling a matrix with both device noise parameters down the diagonal. For the preceding 2-port network this is,
\[
\mathbf{C_t} = 
\renewcommand*{\arraystretch}{1.2}
\begin{bmatrix}
\nwv{c_1}          & \nwcv{c_{1}}{c_{2}}         & 0   & 0   \\
\nwcv{c_{2}}{c_{1}}           & \nwv{c_2}         & 0 &0 \\
0            & 0         & \nwv{c_3} & \nwcv{c_{3}}{c_{4}} \\
0            & 0  & \nwcv{c_{4}}{c_{3}}   & \nwv{c_4}  \\
\end{bmatrix}
\]

Each covariance and variance term can then be computed with,
\[ \nwcv{b_j}{b_k} = \sum{\Big((\mathbf{A}_{j:})^T \mathbf{A}^*_{k:}\Big) \odot \mathbf{C_t}} \]
%
Where \((\mathbf{A}_{j:})^T \mathbf{A}^*_{k:}\) is the outer product of the \(j^{th}\) row of \(\mathbf{A}\) with the conjugate of the \(k^{th}\) row. This forms a matrix the same size as \(\mathbf{C_t}\) and is then multiplied element-wise with \(\mathbf{C_t}\). The result is then summed across all elements to compute the noise wave correlation product \(\nwcv{b_j}{b_k}\). Someone has likely found a more elegant way to compute this, but so far I have been unable to find anything.

Each element of the noise wave parameter matrix \(\mathbf{C_s}\) of the aggregate network are given by \(\nwcv{b_j}{b_k}\). Note that 2 or more ports of the aggregate device are internally connected, and the corresponding rows and columns of \(\mathbf{C_s}\) can be removed to reduce \(\mathbf{C_s}\) to describe only the external ports. The network can then be treated as standalone device and connected with other devices by repeating the process in this section.

One last hurdle is connecting ports together from the same device, as this frequently occurs when building up networks. For example, after the connection in Figure \ref{fig:img/interconnect}, the aggregate 4-port network has its own \(\mathbf{C_s}\) matrix and is now a self-contained device. To connect two ports together from the same device, \(\mathbf{A}\) is composed by inserting columns of the s-matrix from the same device, \(\mathbf{C_t}\) is simply the device noise correlation matrix, and the remaining process is identical.

\subsection{Calculating Noise Figure}
Once the noise wave correlation matrix is determined for a device or network, the noise figure can be calculated by finding the signal to noise ratios at two different ports. The noise figure of a 2-port device is,
\begin{gather}
F = \left(\frac{S_i}{N_1}\right) \left(\frac{N_2}{S_o}\right) = \left(\frac{S_i}{N_1}\right) \left(\frac{S_{21}^2 N_1 +S_{22}^2 N_2+ \nwv{c_2}}{S_{21}^2 S_i}\right) \notag \\
F = \frac{S_{21}^2 N_1 +S_{22}^2 N_2+ \nwv{c_2}}{S_{21}^2 N_1} \label{eq:nf-def} 
\end{gather}
Where \(N_1\) and \(N_2\) are the noise powers entering port 1 and 2, and \(S_i\) is the signal power entering port 1, which is arbitrary since it cancels out.

After some experimenting, it was found that if the noise figure was computed from the equivalent temperature using  (\ref{eq:f2temp}) and (\ref{eq:f2temp0}), instead of using (\ref{eq:nf-def}), the results matched ADS more closely. The problem is apparent by deriving (\ref{eq:f2temp}). First, the device temperature is found by using the condition \(\Gamma_s = 0\), which leads to,
\[T_e = {\nwv{c_2} \over k|S_{21}|^2}\]

The equation to convert \(T_e\) to a noise figure can be derived by taking the signal to noise ratios at the input and output and setting \(N_2 = 0\),

\begin{gather}
F = \left(\frac{S_i}{N_1}\right) \left({S_{21}^2( N_1 + kT_e) \over S_{21}^2 S_i}\right) \notag \\
F =  { N_1 + kT_e \over N_1} = {kT_e \over N_1} + 1\notag
\end{gather}

Since \(N_1 = kT_0\) this reduces to,
\[F = {T_e \over T_0} + 1\]

ADS appears to ignore \(N_2\) by using the temperature method, which usually leads to accurate results for amplifiers because \(N_1\) is usually much larger than the reflected noise wave from port 2 (\(N_2\)). However, this breaks down if \(S_{21}\) is comparable to \(S_{22}\). For most amplifiers, the temperature method will give accurate results. However, the contribution of the reflected noise wave from port 2 (\(N_2\)) should be included for noisy components with low gain.

\vfill

\pagebreak

\vfill

\pagebreak

\begin{thebibliography}{00}
\bibitem{b1} Bertsekas, Dimitri P. \& Tsitsiklis, John N. (2002). Introduction to Probability Vol. 1. Athena Scientific Belmont, Ma.
\bibitem{b2} D. M Pozar, ``Microwave Engineering". Hoboken, NJ, Wiley, 2012
\bibitem{b3} S. W. Wedge and D. B. Rutledge, ``Wave techniques for noise modeling and measurement," in IEEE Transactions on Microwave Theory and Techniques, vol. 40, no. 11, pp. 2004-2012, Nov. 1992, doi: 10.1109/22.168757.
\bibitem{b4} S. W. Wedge and D. B. Rutledge, ``Noise waves and passive linear multiports," in IEEE Microwave and Guided Wave Letters, vol. 1, no. 5, pp. 117-119, May 1991, doi: 10.1109/75.89082.
\bibitem{b5} H. Bosma, ``On the theory of linear noisy systems," Phillips Res. Rep. Suppl., no. 10, 1967
\bibitem{b6} S. Van den Bosch and L. Martens, ``Improved impedance-pattern generation for automatic noise-parameter determination," in IEEE Transactions on Microwave Theory and Techniques, vol. 46, no. 11, pp. 1673-1678, Nov. 1998, doi: 10.1109/22.734556.
\end{thebibliography}



\end{document}





